\section{落地清单：把课程原则转成工程实施步骤}

结合整讲内容，可以给出一个可执行的安全实施路线，适合作为企业 Agent 平台的 baseline：

\begin{table}[H]
\centering
\begin{tabular}{p{1.7cm}p{3.1cm}p{5.8cm}}
\toprule
\textbf{阶段} & \textbf{目标} & \textbf{关键动作} \\
\midrule
S0 资产梳理 & 明确风险边界 & 列出工具、数据、身份、外联目标与高危动作清单 \\
S1 威胁建模 & 定义攻击者与后果 & 按 black/gray/white box 划分能力，明确 attack path 和 impact \\
S2 防御基线 & 建立最小可用防线 & 输入过滤 + 输出约束 + 最小权限 + 人工审批高危动作 \\
S3 执行策略 & 从静态到动态策略 & 引入 DSL 策略执行、上下文动态授权、策略审计日志 \\
S4 评测体系 & 提升可复现性 & 建立统一接口评测、对抗回归测试、持续监测面板 \\
S5 治理闭环 & 安全运营常态化 & 漏洞响应、红队演练、策略迭代、版本回溯与追责机制 \\
\bottomrule
\end{tabular}
\caption{基于本讲内容整理的 Agent 安全工程落地路线}
\end{table}

\begin{knowledgebox}{建议优先落地的三件事}
\begin{enumerate}
    \item 先做 ``权限最小化 + 高危动作审批''，立即降低高影响风险；
    \item 再做 ``统一评测接口 + 对抗回归''，避免盲目上线；
    \item 最后做 ``上下文策略 + 运行时监控''，把系统从可用提升到可控。
\end{enumerate}
\end{knowledgebox}

\subsection{本章小结}
课程思想可以直接映射到工程实践。关键不是一次性搭完所有机制，而是先建立可审计、可演进、可持续监测的安全基座。

